% 结论章节

\section{结论}

本文针对信用卡欺诈检测这一商业智能问题，比较了两种异构深度神经网络架构（FT-Transformer、DeepFM）与XGBoost基线的性能，并提出了Tree-Supervised知识蒸馏方案将XGBoost的树归纳偏置迁移至深度模型。主要发现如下：

\begin{enumerate}
    \item \textbf{DeepFM与XGBoost统计不可区分}：DeepFM（AUPRC=0.9336）达到与XGBoost（0.9331）统计上不可区分的性能水平（单一随机种子下差距仅0.0005）。DeepFM的显式二阶交互在PCA匿名化特征空间中展现出高度鲁棒性。
    \item \textbf{Sparsemax拯救了FT-Transformer}：FT-Transformer经Sparsemax替代Softmax后，AUPRC从0.8738跃升至0.9179（+4.4个百分点），ECE从0.2328降至0.0120（+20倍校准改善），业务成本从269降至99。Sparsemax的稀疏性约束有效缓解了注意力塌缩问题，但注意力分布仍近乎均匀，说明PCA空间本身缺乏供注意力机制利用的跨特征语义关联。
    \item \textbf{TreeReg与FTT性能下降强相关，但因果关系尚未闭合}：FTT+KD消融实验以AUPRC=0.9347（+1.68pp vs. 基线，ECE=0.0007）证明软标签蒸馏对注意力型架构同样高效——推翻了此前"蒸馏是架构依赖的"的预判。FTT从KD获益（+1.68pp）与DeepFM（+1.59pp）高度一致。然而，FTT+KD+TreeReg的断崖式下降（-3.73pp至0.8974）虽与"树先验对嵌入层的刚性约束与Sparsemax注意力拮抗"的假设一致，替代解释——正则化强度$\lambda{=}0.1$可能不当、TreeReg与KD的交互效应、FTT+KD变体与FTT基线的架构参数不完全一致（KD变体使用$\text{d\_model}{=}48$, $\text{n\_layers}{=}2$，基线使用$\text{d\_model}{=}64$, $\text{n\_layers}{=}3$）——尚未通过控制实验排除。因此，当前证据支持KD在本研究条件下对两种架构均安全有效，但TreeReg的精确因果角色需未来多种子、多$\lambda$扫描实验加以确认。软标签蒸馏是可靠的性能提升手段；嵌入层正则化在注意力型架构上的使用需极其谨慎。
    \item \textbf{商业成本矩阵至关重要}：FTT+KD实现最低业务成本（97），而AUPRC领先的DeepFM+KD成本为120，XGBoost为215，揭示了精度与成本间的非单调关系，为业务决策提供帕累托前沿分析依据。
\end{enumerate}

综上，深度表格模型在信用卡欺诈检测领域已展现出与XGBoost匹敌的潜力——DeepFM+KD以AUPRC=0.9495确立最优性能，FTT+KD以0.9347和ECE=0.0007实现最佳校准与最低成本（Cost=97）。Tree-Supervised蒸馏的消融实验在本研究条件下表明：软标签蒸馏（3A）对两种架构均安全有效（FTT+1.68pp, DeepFM+1.59pp），而TreeReg（3B）与FTT性能的大幅下降强相关（-3.73pp），其精确因果机制（结构性冲突 vs. 强度不当 vs. 交互效应）是未来研究的明确方向。需要指出，FTT基线（$\text{d\_model}{=}64$, $\text{n\_layers}{=}3$）与KD变体（$\text{d\_model}{=}48$, $\text{n\_layers}{=}2$）的架构参数不完全一致，KD消融实验的效应量可能受到混淆——这是本研究最核心的方法论局限，应在解读所有蒸馏相关结论时予以考量。未来工作应聚焦于：统一架构参数的控制实验、多种子统计检验与KD超参敏感性分析、为注意力模型设计去冲突的树先验注入方式。

本研究为商业智能课程的实践学习提供了完整的端到端案例：从数据集获取与预处理、多个深度学习模型的实现与训练、多维度的评估分析，到基于可解释性的商业决策支持。
